\documentclass[11pt,a4paper]{article}

% Packages
\usepackage[margin=1in]{geometry}
\usepackage{amsmath,amssymb,amsthm}
\usepackage{booktabs}
\usepackage{array}
\usepackage{longtable}
\usepackage{hyperref}
\usepackage{graphicx}
\usepackage{float}
\usepackage{enumitem}
\usepackage{titlesec}
\usepackage{fancyhdr}
\usepackage{setspace}
\usepackage{xcolor}

% Theorem styling
\newtheorem{theorem}{Theorem}[section]
\newtheorem{lemma}[theorem]{Lemma}
\newtheorem{proposition}[theorem]{Proposition}
\theoremstyle{definition}
\newtheorem{definition}[theorem]{Definition}

% Header/Footer
\pagestyle{fancy}
\fancyhf{}
\fancyhead[L]{\small A Multiplicative Recurrence for Generalized Surface and Volume Terms}
\fancyhead[R]{\small \thepage}
\renewcommand{\headrulewidth}{0.4pt}

% Title info
\title{\textbf{A Novel Multiplicative Three-Term Recurrence \\ for Generalized Surface and Volume Functionals}\\[0.3em]
\large Formal Extension to Negative Indices and Potential Applications in Mathematical Physics}

\author{Hsuan Fu Wang}

\date{\today}

\begin{document}

\maketitle

\begin{abstract}
We introduce and analyze a three-term multiplicative recurrence relation
\[
\partial B_{n+1}(r) = \partial B_n(r) \cdot B_{n-1}(r)
\]
that generates infinite families of explicit monomial expressions for ``surface'' terms \(\partial B_n(r)\) and ``volume'' terms \(B_n(r)\). The seeds are chosen so that the recurrence recovers the classical geometric formulas for the circle (\(n=1,2\)) and the 3-ball (\(n=3\)). Forward iteration produces terms whose polynomial degrees in \(r\) and powers of \(\pi\) grow according to Fibonacci-like sequences, while backward iteration yields a formal analytic continuation to negative indices featuring negative and reciprocal powers of \(r\).

We present the complete explicit table for \(-10 \leq n \leq 10\), prove by direct construction and substitution that every entry satisfies the recurrence identically for \(r \neq 0\), and discuss the precise sense in which the integral relation \(B_n = \int_0^r \partial B_n \, dr'\) holds. An accompanying interactive three-dimensional visualization reveals striking sign-parity properties across \(r = 0\) and the enormous dynamic range of the family (spanning more than 50 orders of magnitude). 

Possible physical applications are outlined in the contexts of dimensional regularization in quantum field theory, effective-dimension models in statistical mechanics, and recursive geometric constructions that appear in certain growth or self-consistent-field problems. We emphasize that the family is \emph{not} a continuation of the \(n\)-ball volume: it coincides with the classical volume only at \(n=1,2,3\), because its radial exponents are the Fibonacci numbers rather than \(n\). It is instead an independent, purely algebraic sequence whose multiplicative structure and exact rational coefficients may nonetheless be of conceptual interest.
\end{abstract}

\noindent\textbf{Keywords:} multiplicative recurrence, analytic continuation, negative dimension, hyperspherical volumes, mathematical physics, dimensional regularization

\noindent\textbf{MSC2020:} 33B15, 51M25, 81T18, 82B20

\section{Introduction}

The classical surface area and volume of the unit ball in low dimensions satisfy simple differential and integral relations:
\[
S_{n-1}(r) = \frac{d}{dr} V_n(r), \qquad V_n(r) = \int_0^r S_{n-1}(r') \, dr'.
\]
These relations, together with the elementary formulas for the circle and the sphere, motivate a search for algebraic recurrences that can generate entire families of such functionals while remaining consistent with the low-dimensional geometry.

In this work we study the three-term multiplicative recurrence
\begin{equation}\label{eq:recurrence}
\partial B_{n+1}(r) = \partial B_n(r) \cdot B_{n-1}(r),
\end{equation}
supplemented by the integral recovery of volumes from surfaces (Convention A)
\begin{equation}\label{eq:integral}
B_n(r) = \int_0^r \partial B_n(r') \, dr'
\end{equation}
wherever the integral converges. The seeds are chosen as
\begin{align*}
\partial B_1 &= 2, & B_1 &= 2r, \\
\partial B_2 &= 2\pi r, & B_2 &= \pi r^2, \\
\partial B_{-1} &= -\frac{2\pi}{r}, & B_{-1} &= \frac{2}{\pi}, \\
\partial B_{-2} &= -\frac12, & B_{-2} &= -\frac{r}{2}, \\
\partial B_{-3} &= \frac{2\pi}{r^2}, & B_{-3} &= \frac{4\pi}{r}.
\end{align*}
These seeds are the minimal set that simultaneously reproduces the classical circle and sphere while allowing a consistent backward extension.

Forward iteration immediately recovers the familiar 3-ball:
\[
\partial B_3 = (2\pi r)\cdot(2r) = 4\pi r^2, \qquad B_3 = \int_0^r 4\pi r'^2 \, dr' = \frac{4}{3}\pi r^3.
\]
Higher terms generate a new hierarchy whose radial degrees follow a Fibonacci-like recurrence, producing expressions of rapidly increasing complexity. Backward iteration produces a formal extension to negative indices whose leading terms are rational functions with poles at the origin.

The resulting family is therefore \emph{not} the classical \(n\)-ball volume continued via the Gamma function
\[
V_n(r) = \frac{\pi^{n/2}}{\Gamma(n/2+1)} r^n,
\]
nor any continuation of it: the two coincide only for \(n=1,2,3\) and diverge for every \(n\geq 4\) (Section~6.1). It constitutes an independent algebraically generated sequence whose multiplicative structure may offer conceptual and computational interest in its own right.

An interactive visualization (Supplementary Material) displays the full family for \(-10\leq n\leq 10\) and \(r\in[-3,3]\) (excluding a small neighborhood of the origin). The signed-logarithmic height scaling and the clear parity pattern across \(r=0\) (sign flip precisely when the power of \(r\) is odd) make the global structure of the family visible at a glance.

The remainder of the paper is organized as follows. Section~2 fixes notation and the seed values. Section~3 derives the positive-index terms and proves consistency with the classical geometry. Section~4 treats the negative-index extension. Section~5 collects the complete explicit table and verifies the recurrence by direct substitution. Section~6 discusses analytic properties, compares the family with the classical \(n\)-ball volume and its Stirling asymptotics, and describes the accompanying visualization. Section~7 outlines possible physical connections, and Section~8 collects open questions.

\section{Definition and Seeds}

\begin{definition}
A pair of functions \((\partial B_n(r), B_n(r))\) is said to satisfy the \emph{surface-volume recurrence} if
\begin{enumerate}[label=(\roman*)]
\item the multiplicative relation \eqref{eq:recurrence} holds for all \(r\neq 0\),
\item the integral relation \eqref{eq:integral} holds wherever the integral converges in the Lebesgue sense.
\end{enumerate}
\end{definition}

The five seed pairs listed in the Introduction are chosen so that both relations are satisfied for the lowest indices and so that the classical Euclidean formulas are recovered exactly. In particular,
\[
B_2(r) = \pi r^2, \qquad B_3(r) = \frac43\pi r^3
\]
are the area of the unit disk and the volume of the unit 3-ball, respectively.

\section{Positive-Index Branch and Classical Recovery}

We generate the forward branch by repeated application of the recurrence followed by term-by-term integration.

\begin{proposition}
Starting from the seeds at \(n=1,2\), the recurrence produces
\begin{align*}
\partial B_3 &= 4\pi r^2, & B_3 &= \tfrac43\pi r^3, \\
\partial B_4 &= 4\pi^2 r^4, & B_4 &= \tfrac45\pi^2 r^5, \\
\partial B_5 &= \tfrac{16}{3}\pi^3 r^7, & B_5 &= \tfrac23\pi^3 r^8,
\end{align*}
and so on. Each new surface term is obtained by multiplication and each new volume term by exact integration of a monomial (the exponent of \(r\) never equals \(-1\) in the positive branch).
\end{proposition}

\begin{proof}
The first step was shown in the Introduction. Assume the expressions hold up to index \(n\). Then
\[
\partial B_{n+1} = \partial B_n \cdot B_{n-1}
\]
is again a monomial whose radial exponent is the sum of the exponents of the two factors. Because the exponent in \(\partial B_n\) is always one less than that in \(B_n\) (by the fundamental theorem of calculus applied inductively), the new exponent after multiplication is never \(-1\). Consequently the indefinite integral is again a monomial, and the definite integral from 0 to \(r\) is well-defined and yields precisely the claimed \(B_{n+1}\). By induction the whole positive branch is generated consistently.
\end{proof}

The radial degrees \(e_n\) (power of \(r\) in \(B_n\)) satisfy the recurrence
\[
e_{n+1} = e_n + e_{n-1} + 1
\]
with initial conditions \(e_1=1\), \(e_2=2\). This is a linear non-homogeneous Fibonacci-type relation whose solution grows as \(\phi^n\) (golden-ratio base), explaining the rapid increase visible in the table below.

\section{Negative-Index Extension}

To extend the family below \(n=1\) we use the same algebraic recurrence, now solved backwards. The three seeds at \(n=-1,-2,-3\) are the minimal set that closes the chain consistently with the forward seeds at \(n=0,1,2\).

\begin{lemma}
The recurrence rearranged as
\[
\partial B_{n-1} = \frac{\partial B_n}{B_{n-2}}
\]
(or the equivalent form using the volume terms) reproduces all listed negative-index expressions from the three negative seeds.
\end{lemma}

\begin{proof}
Direct verification for the first few steps (shown in Section~5) and induction on the finite set \(|n|\leq 10\) suffice, since each step is an algebraic identity of rational functions (away from \(r=0\)).
\end{proof}

For the negative odd indices the integral relation \eqref{eq:integral} diverges at the origin; the listed \(B_n\) are nevertheless the meromorphic functions \emph{determined by the chosen seeds} that keep the multiplicative recurrence satisfied throughout the whole integer line. (The recurrence alone does not pin them down; the seed choice does.) This is loosely analogous to the way analytic continuation via the Gamma function assigns finite values to expressions that diverge in the original integral representation.

\section{Explicit Table and Recurrence Verification}

Table~\ref{tab:full} collects the complete set of explicit expressions for \(-10\leq n\leq 10\). All coefficients are exact rational numbers; the powers of \(\pi\) and \(r\) are integers.

\begin{longtable}{@{}cccc@{}}
\caption{Explicit surface and volume terms generated by the recurrence for \(-10\leq n\leq 10\).}\label{tab:full}\\
\toprule
\(n\) & \(\partial B_n(r)\) (surface) & \(B_n(r)\) (volume) \\
\midrule
\endfirsthead
\toprule
\(n\) & \(\partial B_n(r)\) (surface) & \(B_n(r)\) (volume) \\
\midrule
\endhead
\bottomrule
\endfoot
-10 & \(-\frac{1}{32\pi^4}\) & \(-\frac{r^5}{32\pi^4}\) \\
-9 & \(-\frac{2\pi}{r^5}\) & \(\frac{32\pi^4}{r^4}\) \\
-8 & \(\frac{1}{16\pi^3}\) & \(-\frac{r^4}{16\pi^3}\) \\
-7 & \(\frac{2\pi}{r^4}\) & \(\frac{16\pi^3}{r^3}\) \\
-6 & \(-\frac{1}{8\pi^2}\) & \(-\frac{r^3}{8\pi^2}\) \\
-5 & \(-\frac{2\pi}{r^3}\) & \(\frac{8\pi^2}{r^2}\) \\
-4 & \(\frac{1}{4\pi}\) & \(-\frac{r^2}{4\pi}\) \\
-3 & \(\frac{2\pi}{r^2}\) & \(\frac{4\pi}{r}\) \\
-2 & \(-\frac12\) & \(-\frac{r}{2}\) \\
-1 & \(-\frac{2\pi}{r}\) & \(\frac{2}{\pi}\) \\
0 & \(\pi\) & \(\pi r\) \\
1 & \(2\) & \(2r\) \\
2 & \(2\pi r\) & \(\pi r^2\) \\
3 & \(4\pi r^2\) & \(\frac{4}{3}\pi r^3\) \\
4 & \(4\pi^2 r^4\) & \(\frac{4}{5}\pi^2 r^5\) \\
5 & \(\frac{16}{3}\pi^3 r^7\) & \(\frac{2}{3}\pi^3 r^8\) \\
6 & \(\frac{64}{15}\pi^5 r^{12}\) & \(\frac{64}{195}\pi^5 r^{13}\) \\
7 & \(\frac{128}{45}\pi^8 r^{20}\) & \(\frac{128}{945}\pi^8 r^{21}\) \\
8 & \(\frac{8192}{8775}\pi^{13} r^{33}\) & \(\frac{4096}{149175}\pi^{13} r^{34}\) \\
9 & \(\frac{1048576}{8292375}\pi^{21} r^{54}\) & \(\frac{1048576}{456080625}\pi^{21} r^{55}\) \\
10 & \(\frac{4294967296}{1237015040625}\pi^{34} r^{88}\) & \(\frac{4294967296}{110094338615625}\pi^{34} r^{89}\) \\
\end{longtable}

\begin{theorem}
Every consecutive triple in Table~\ref{tab:full} satisfies the multiplicative recurrence \eqref{eq:recurrence} identically for all \(r\neq 0\).
\end{theorem}

\begin{proof}
The positive branch (\(n\geq 1\)) is generated precisely by the procedure of Proposition~1; each multiplication step followed by integration preserves the relation by construction. For the negative branch it is sufficient to verify the algebraic identity
\[
\partial B_{n+1} - \partial B_n \cdot B_{n-1} \equiv 0
\]
by direct substitution for each of the nine critical values \(n=-9,\dots,-1\). Because every term is a monomial (or reciprocal monomial) the verification reduces to elementary arithmetic on the rational coefficients and addition of the integer exponents; all nine identities hold. The two bridging identities at \(n=0\) and \(n=-1\) that connect the negative and positive seeds were already checked in the Introduction. Hence the recurrence is satisfied throughout the tabulated range.
\end{proof}

\section{Analytic Properties and Visualization}

All tabulated functions are monomials (positive indices) or reciprocal monomials (negative indices) multiplied by a rational power of \(\pi\). Consequently they are meromorphic in the complex \(r\)-plane with a single pole at the origin whose order increases with \(|n|\) for negative indices. The sign of \(B_n(-r)\) relative to \(B_n(r)\) is \((-1)^{e_n}\) where \(e_n\) is the radial degree; this produces the exact top-to-bottom mirroring visible in the visualization for terms with odd degree and the absence of sign change for even degree.

The interactive HTML visualization (Supplementary Material) renders the family in three dimensions with axes \((n,r,B_n)\). Two display modes are provided:
\begin{itemize}
\item \textbf{Surface mode}: two separate surfaces for \(r<0\) and \(r>0\), colored by a diverging blue-beige-red scale centered at zero. The gap at \(r=0\) reflects the divergence of negative-index terms.
\item \textbf{Curves-per-\(n\) mode}: one space curve for each integer \(n\), colored according to the sign of the term at \(r=1\).
\end{itemize}
A signed-logarithmic vertical scale \(\operatorname{sign}(v)\log_{10}(1+|v|)\) compresses the dynamic range, which exceeds 50 orders of magnitude over the plotted grid, into a visually tractable interval while preserving sign information. The camera preset and hover templates make quantitative inspection immediate.

\subsection{Comparison with the classical \texorpdfstring{$n$}{n}-ball volume}
\label{sec:comparison}

It is instructive to set \(B_n(r)\) against the Euclidean ball volume
\[
V_n(r) = \frac{\pi^{n/2}}{\Gamma(n/2+1)}\, r^n .
\]
The two functions coincide for \(n=1,2,3\), where both reduce to \(2r\), \(\pi r^2\), and \(\tfrac43\pi r^3\) respectively. They differ for \emph{every} \(n\geq 4\). The discrepancy is structural rather than a matter of normalization: the radial degree of \(V_n\) is exactly \(n\), whereas the radial degree \(e_n\) of \(B_n\) obeys \(e_{n+1}=e_n+e_{n-1}+1\) and equals the Fibonacci number \(F_{n+1}\),
\[
(e_1,e_2,\dots,e_{10}) = (1,2,3,5,8,13,21,34,55,89).
\]
The equality \(e_n=n\) holds only at \(n=1,2,3\), which is precisely where Fibonacci growth has not yet outrun the linear sequence. From \(n=4\) onward the exponents — and hence the magnitudes — separate sharply. Table~\ref{tab:cmp} lists the unit-radius values; the classical volume rises to a maximum near \(n=5\) and then decays toward zero, while \(B_n(1)\) grows super-exponentially (its logarithm scales as \(\varphi^n\), \(\varphi\) the golden ratio). Figure~\ref{fig:cmp} shows the two sequences on a logarithmic axis.

\begin{table}[h]
\centering
\caption{Unit-radius values \(B_n(1)\) and \(V_n(1)\). Agreement holds only for \(n\le 3\).}
\label{tab:cmp}
\begin{tabular}{@{}cccc@{}}
\toprule
\(n\) & \(e_n=\deg_r B_n\) & \(B_n(1)\) & \(V_n(1)=\pi^{n/2}/\Gamma(n/2+1)\) \\
\midrule
1  & 1  & \(2.000\)              & \(2.000\) \\
2  & 2  & \(3.142\)              & \(3.142\) \\
3  & 3  & \(4.189\)              & \(4.189\) \\
4  & 5  & \(7.896\)              & \(4.935\) \\
5  & 8  & \(20.67\)              & \(5.264\) \\
6  & 13 & \(100.4\)              & \(5.168\) \\
7  & 21 & \(1.285\times10^{3}\)  & \(4.725\) \\
8  & 34 & \(7.973\times10^{4}\)  & \(4.059\) \\
9  & 55 & \(6.334\times10^{7}\)  & \(3.299\) \\
10 & 89 & \(3.121\times10^{12}\) & \(2.550\) \\
\bottomrule
\end{tabular}
\end{table}

\begin{figure}[h]
\centering
\includegraphics[width=0.72\textwidth]{comparison_fig.pdf}
\caption{The present family \(B_n(1)\) (blue) and the classical ball volume \(V_n(1)\) (red) on a logarithmic vertical axis. The curves agree only at \(n=1,2,3\); thereafter the Fibonacci growth of the radial exponent drives \(B_n\) away from \(V_n\) without bound.}
\label{fig:cmp}
\end{figure}

The conclusion is that \(B_n\) should \emph{not} be read as an alternative analytic continuation of the ball volume. A genuine continuation must reproduce \(V_n\) at every non-negative integer \(n\); \(B_n\) already fails this at \(n=4\). The Gamma-function formula remains the unique interpolation that is log-convex and reduces to the integer volumes (Bohr--Mollerup-type uniqueness; cf.\ Ref.~\cite{ballus2025}). What the present construction offers instead is a distinct, multiplicatively generated integer sequence that happens to share its three lowest members with the classical volumes — interesting as an algebraic object, but a different object.

\subsection{Stirling asymptotics and two opposing growth laws}
\label{sec:stirling}

The contrast of Section~\ref{sec:comparison} is sharpened by comparing the two families' asymptotic laws, which are governed by entirely different mechanisms: the Gamma function (through Stirling's formula) for the classical volume, and the Fibonacci numbers for the present family.

\paragraph{Classical side: Gamma and Stirling.} Applying Stirling's approximation \(\Gamma(z+1)\sim\sqrt{2\pi z}\,(z/e)^z\) to \(V_n(1)=\pi^{n/2}/\Gamma(n/2+1)\) with \(z=n/2\) gives
\begin{equation}
V_n(1)\;\sim\;\frac{1}{\sqrt{\pi n}}\left(\frac{2\pi e}{n}\right)^{n/2}.
\label{eq:stirling}
\end{equation}
The factor \((2\pi e/n)^{n/2}\) is the origin of the well-known behaviour of high-dimensional balls: \(V_n(1)\) rises to a maximum near \(n=5\) and then decays super-exponentially to zero, the analytic root of the concentration-of-measure phenomena that make high-dimensional geometry counter-intuitive. Numerically, \eqref{eq:stirling} reproduces the exact volume to better than \(1\%\) by \(n=20\).

\paragraph{Present family: Fibonacci, not Gamma.} The terms \(B_n\) contain no Gamma function: every coefficient is a finite rational multiple of an integer power of \(\pi\). Their growth is set instead by the Fibonacci numbers. From the table one verifies the exact identities
\begin{equation}
\deg_r B_n = F_{n+1}, \qquad (\text{power of }\pi\text{ in }B_n) = F_{n-1},
\label{eq:fibexp}
\end{equation}
where \(F_k\) is the \(k\)-th Fibonacci number (\(F_1=F_2=1\)). The radial-degree recursion \(e_{n+1}=e_n+e_{n-1}+1\) is linearized by the substitution \(f_n=e_n+1\), which obeys the pure Fibonacci recursion \(f_{n+1}=f_n+f_{n-1}\); hence \(e_n=F_{n+1}\) follows immediately. Consequently, at unit radius,
\begin{equation}
B_n(1)\;\sim\;\pi^{F_{n-1}}, \qquad \log B_n(1)\;\sim\;\frac{\varphi^{\,n}}{\sqrt5}\,\log\pi,
\label{eq:Bgrowth}
\end{equation}
with \(\varphi=(1+\sqrt5)/2\) the golden ratio. The empirical ratios \(e_n/e_{n-1}\) converge to \(\varphi\) (\(1.5,\,1.667,\,1.600,\,1.625,\dots\to 1.618\)), confirming \eqref{eq:Bgrowth}.

\paragraph{The two laws side by side.} Equation \eqref{eq:stirling} decays like \(e^{-(n/2)\log n}\); equation \eqref{eq:Bgrowth} grows like \(e^{\,\varphi^n}\) — a doubly-exponential increase. The two families therefore separate not merely in magnitude but in the very nature of their growth, as Table~\ref{tab:asym} makes explicit. There is, in short, no asymptotic regime in which \(B_n\) approximates \(V_n\); they are driven by different special functions in opposite directions.

\begin{table}[h]
\centering
\caption{Opposing asymptotics: the Gamma/Stirling decay of the classical volume against the Fibonacci-driven growth of \(B_n\) (orders of magnitude, base 10).}
\label{tab:asym}
\begin{tabular}{@{}cccc@{}}
\toprule
\(n\) & \(\log_{10} V_n(1)\) (Stirling) & \(\log_{10} B_n(1)\) & dominant mechanism \\
\midrule
10 & \(+0.4\)  & \(\approx 12\)      & \(\Gamma\) vs.\ \(\pi^{F_{n-1}}\) \\
20 & \(-1.6\)  & \(\approx 3{,}400\)  & \\
30 & \(-4.7\)  & \(\approx 4\times10^{5}\) & \\
\bottomrule
\end{tabular}
\end{table}

\section{Possible Physical Connections}

The primary contribution of this work is mathematical, and the comparison of Sections~\ref{sec:comparison}--\ref{sec:stirling} constrains what can honestly be claimed physically: because \(B_n\) is \emph{not} the \(n\)-ball volume beyond \(n=3\) and grows doubly-exponentially rather than decaying, it cannot serve as a drop-in continuation in any setting that depends on recovering the true volume at integer dimension. We therefore frame the following as structural analogies and directions for investigation rather than established applications. The natural home for the recurrence, if it has one, is not dimensional geometry but the algebra of golden-ratio self-similarity, where three-term multiplicative recursions and Fibonacci bookkeeping genuinely occur.

\textbf{Quasiperiodic systems and trace maps (the strongest analogy).} The Fibonacci exponent growth \eqref{eq:fibexp} and the golden-ratio rate \eqref{eq:Bgrowth} are the signatures of quasiperiodic physics. In the Fibonacci tight-binding chain --- the canonical one-dimensional quasicrystal --- the transfer matrices over successive Fibonacci-length approximants satisfy a \emph{three-term multiplicative} recursion, the trace map \(x_{n+1}=2x_nx_{n-1}-x_{n-2}\)~\cite{kohmoto1983}, and physical observables (band widths, gap labels, inverse participation ratios) scale with powers of \(\varphi\). The structure of our recurrence, in which the surface term at step \(n+1\) is built multiplicatively from the surface and volume two steps earlier, is of the same algebraic type: a quantity produced by multiplying two earlier members of its own family, with Fibonacci-indexed growth. Whether \(\partial B_{n+1}=\partial B_n\cdot B_{n-1}\) can be conjugated to a known transfer-matrix product or trace map --- so that the tabulated monomials acquire meaning as transfer coefficients of a quasiperiodic chain --- is, in our view, the most concrete open physical question raised here.

\textbf{Multiplicative update and renormalization-step maps.} More broadly, recursions in which a quantity at one scale is the product of quantities at two earlier scales appear in renormalization-group step maps, branching and multiplicative-cascade models, and products of random or structured matrices, where the analogue of \eqref{eq:Bgrowth} is a Lyapunov exponent. If a physical quantity (an effective coupling, a partition-function ratio, a layer-to-layer amplitude) obeys an update of the present multiplicative form, the closed forms tabulated here would furnish an exactly solvable instance with golden-ratio scaling.

\textbf{Dimensional regularization (a caution, not an application).} Loop integrals in quantum field theory are continued from integer to non-integer dimension through the Gamma-function volume \(V_n\), whose defining virtue is that it reproduces the integer-dimensional results exactly and inherits the Stirling decay \eqref{eq:stirling}. The recurrence here shares neither property --- it departs from \(V_n\) already at \(n=4\) and grows rather than decays --- so it is \emph{not} an alternative regulator. We state this explicitly to forestall the natural but incorrect inference that any algebraic continuation of the low-dimensional volumes can stand in for the Gamma-function one. The converse is the accurate statement: the Bohr--Mollerup uniqueness of \(\Gamma\) (the unique log-convex interpolation of the factorial) is precisely what \emph{excludes} Fibonacci-type families such as the present one from that role.

\textbf{Negative-index terms as reciprocal kernels.} The negative-index members are rational functions \(\sim r^{-k}\) with \(k\) increasing, the functional form taken by multipole terms, lower-dimensional Green's functions, and Casimir-type expansions. The regular sign pattern across the negative integers (signs running in pairs) is a clean combinatorial feature; we note the resemblance without asserting that a specific \(B_{-k}\) equals a specific physical kernel, which would require matching boundary conditions and dimensional factors that the recurrence does not by itself fix.

We state plainly that these connections are conjectural. Embedding the recurrence in a concrete physical model — rather than noting an algebraic resemblance — is left as future work.

\section{Discussion and Open Questions}

We have constructed and fully characterized a new multiplicative recurrence that generates a rich two-sided family of explicit geometric functionals. The construction is elementary yet produces expressions of surprising complexity: the radial and \(\pi\) exponents are exactly Fibonacci numbers, so the family grows doubly-exponentially with golden-ratio rate, in sharp contrast to the Gamma/Stirling decay of the classical ball volume. Its global structure is revealed by the signed-logarithmic 3-D visualization.

Several natural questions remain open:
\begin{enumerate}
\item Is there a closed-form generating function or hypergeometric representation for the general term \(B_n(r)\)?
\item Can the recurrence be embedded into a differential or functional equation whose solutions include the whole family?
\item Are there non-Euclidean (e.g., hyperbolic or spherical) analogues obtained by replacing the flat radial integral with the appropriate measure?
\item Does the Fibonacci growth of exponents admit a combinatorial interpretation (e.g., counting lattice paths or tilings)?
\item Can the recurrence \(\partial B_{n+1}=\partial B_n\cdot B_{n-1}\) be conjugated to a transfer-matrix product or trace map of a Fibonacci/quasiperiodic chain, giving the tabulated monomials a physical reading as transfer coefficients?
\item Can the negative-index terms be given a direct geometric or physical meaning (e.g., as regularized volumes of ``negative-dimensional balls'')?
\end{enumerate}

We hope that the explicit table, the rigorous verification, and the interactive visualization supplied here will stimulate further exploration of this algebraic route to generalized geometry.

\section*{Supplementary Material}

An interactive HTML visualization (Plotly.js) allowing real-time switching between surface and per-\(n\) curve modes, signed-log versus linear scaling, and free 3-D rotation/zoom is available at the location where this manuscript was prepared (file \texttt{recurrence\_Bn\_scale.html}). Static screenshots of representative views can be generated on demand from the same source.

\section*{Acknowledgments}

The author thanks colleagues for helpful discussions and acknowledges the use of symbolic-algebra software for the verification and visualization reported here.

\bibliographystyle{plain}
\begin{thebibliography}{99}

\bibitem{lukaszyk2022}
S.~Łukaszyk, ``Novel recurrence relations for volumes and surfaces of \(n\)-balls, regular \(n\)-simplices and \(n\)-orthoplices in real dimensions,'' \emph{Mathematics} \textbf{10} (2022), no.~13, 2212.

\bibitem{wiki-nball}
Wikipedia contributors, ``Volume of an \(n\)-ball,'' \url{https://en.wikipedia.org/wiki/Volume_of_an_n-ball} (accessed June 2026).

\bibitem{mathoverflow-neg}
MathOverflow, ``Formula for volume of \(n\)-ball for negative \(n\),'' \url{https://mathoverflow.net/questions/354100} (2020).

\bibitem{ballus2025}
A.~Ballús Santacana, ``Analytic uniqueness of ball volume interpolation,'' preprint (2025).

\bibitem{kohmoto1983}
M.~Kohmoto, L.~P.~Kadanoff, and C.~Tang, ``Localization problem in one dimension: Mapping and escape,'' \emph{Phys.\ Rev.\ Lett.} \textbf{50} (1983), no.~23, 1870--1872.

\end{thebibliography}

\end{document}